\documentclass[12pt]{article}
\PassOptionsToPackage{hyphens}{url}
\usepackage{sbc-template}
\usepackage{graphicx,url}
\usepackage[utf8]{inputenc}
\usepackage[brazil]{babel}
\usepackage{amsmath}
\usepackage{float}
\sloppy
\raggedbottom
\title{Manchas Criminais em Alagoas: Uma Análise Espaço-Temporal dos\\
Crimes Violentos Letais Intencionais para Apoio à Tomada de\\
Decisão em Segurança Pública}
\author{Antonio Guilherme da Silva\inst{1}, Antonio Rodrigo L. A. Tenório\inst{1},
  Sandro Gomes Paulino\inst{1}}
\address{Instituto de Computação -- Universidade Federal de Alagoas (UFAL)\\
  Maceió -- AL -- Brasil
  \email{ags2@ic.ufal.br, arlat@ic.ufal.br, sgp@ic.ufal.br}
}
\begin{document}
\maketitle
\begin{abstract}
  The state of Alagoas historically presents one of the highest rates of
  intentional lethal violent crimes (CVLI) in Brazil. This paper presents a
  spatio-temporal analysis of a public CVLI microdata set covering 2012--2026
  (21,217 raw records; 20,369 after cleaning). Through exploratory analysis,
  temporal heatmaps, and geographic visualisation, we identified that 60.36\%
  of crimes occur during nighttime and early morning (18h--5h), 37.88\% on
  weekends, and Maceió accounts for 33.33\% of all state occurrences. As an
  operational component, an XGBoost Poisson regressor was trained on a full
  scenario grid, normalised by demographic population estimates from the 2022
  IBGE Census, to score relative risk by profile. A comparison of supervised
  classifiers (Logistic Regression, Random Forest and XGBoost) was also carried
  out; the XGBoost classifier reached 92.42\% accuracy and 94.70\% ± 0.71\% in
  5-fold cross-validation, but these figures are optimistic and not reliable,
  because the risk label is derived from whole-sample aggregates and causes
  target leakage (Section~6). Results support a data-driven predictive
  patrolling protocol for dynamic resource allocation, subject to the caveats
  discussed.
\end{abstract}
\begin{resumo}
  Alagoas historicamente apresenta um dos maiores índices de Crimes Violentos
  Letais Intencionais (CVLI) do Brasil. Este trabalho apresenta uma análise
  espaço-temporal de uma base pública de microdados de CVLI (2012--2026):
  21.217 registros brutos e 20.369 após a limpeza. Por meio de análise
  exploratória, mapas de calor e visualização geográfica, identificou-se que
  60,36\% dos crimes ocorrem no período noturno e de madrugada (18h--5h),
  37,88\% nos fins de semana, e Maceió responde por 33,33\% das ocorrências.
  Como componente operacional, treinou-se um XGBoost Poisson Regressor sobre
  uma matriz completa de cenários, normalizado por estimativas populacionais
  do Censo IBGE 2022, para pontuar o risco relativo por perfil. Também foi
  realizada uma comparação de classificadores supervisionados (Regressão
  Logística, Random Forest e XGBoost); o classificador XGBoost atingiu 92,42\%
  de acurácia e 94,70\% ± 0,71\% em validação cruzada de 5 \textit{folds}, mas
  esses valores são otimistas e não confiáveis, pois o rótulo de risco é
  derivado de agregados de toda a base, gerando vazamento de alvo (Seção~6).
  Os resultados embasam um protocolo de patrulhamento preditivo orientado a
  dados, sujeito às ressalvas discutidas.
\end{resumo}
\section{Introdução}
O Brasil é um dos países com maior número absoluto de homicídios no mundo,
e o estado de Alagoas historicamente figura entre as unidades federativas com
maiores taxas de mortalidade violenta por habitante~\cite{ipea:2021}. Os
Crimes Violentos Letais Intencionais (CVLI) — categoria que engloba homicídios
dolosos, latrocínios e lesões corporais seguidas de morte — representam não
apenas uma tragédia humana de proporções alarmantes, mas também um grave
problema socioeconômico, com impactos diretos sobre o desenvolvimento
regional, o capital social e a qualidade de vida da
população~\cite{cerqueira:2014}.

Apesar da gravidade do problema, os recursos de segurança pública são
frequentemente alocados de maneira uniforme e reativa, sem considerar os
padrões espaço-temporais da criminalidade. Abordagens tradicionais de
policiamento tendem a responder ao crime após sua ocorrência, em vez de
antecipá-lo por meio da análise de dados históricos~\cite{sherman:1989}.
Com o advento da Ciência de Dados e da Aprendizagem de Máquina, tornou-se
tecnicamente viável transformar grandes volumes de registros históricos de
criminalidade em inteligência operacional capaz de orientar decisões de
policiamento preventivo~\cite{buczak:2016}.

O presente trabalho aborda a seguinte questão de pesquisa: \textit{Existem
padrões de sazonalidade — por dia da semana, mês do ano ou faixa de horário
— associados a surtos de violência em microrregiões específicas de Alagoas
que precedem aumentos nos índices de CVLI?} Para respondê-la, foi realizada
uma análise espaço-temporal sobre a base pública de microdados de CVLI
disponibilizada pela Secretaria de Estado da Segurança Pública de Alagoas
(SSP-AL), cobrindo 14 anos de registros (2012--2026).

As principais contribuições deste trabalho são: (i)~identificação de padrões
temporais e geográficos de concentração do CVLI em Alagoas; (ii)~construção
e avaliação de um modelo preditivo de risco relativo baseado em XGBoost
Poisson, normalizado por população exposta; e (iii)~proposição de
recomendações acionáveis para gestores de segurança pública fundamentadas
nos dados analisados.

O restante deste artigo está organizado da seguinte forma: a Seção~2
descreve a base de dados e o pré-processamento; a Seção~3 apresenta a
análise exploratória; a Seção~4 detalha os métodos aplicados; a Seção~5
apresenta as métricas e a estratégia de avaliação; a Seção~6 discute os
resultados; e a Seção~7 conclui o trabalho.

\section{Base de Dados}
\subsection{Origem e Descrição}
A base de dados utilizada é a \textit{CVLI~AL -- Base de Microdados},
disponibilizada publicamente pelo Portal de Dados Abertos do Estado de
Alagoas, mantido pela Secretaria de Estado da Segurança Pública
(SSP-AL)~\cite{sspAL:2024}. O conjunto de dados abrange o período de
janeiro de 2012 a abril de 2026, totalizando \textbf{21.217 registros} e
\textbf{14 colunas}, entre as quais: identificador do registro
(\texttt{ID\_CONTROLE}), tipo de crime (\texttt{SUBJETIVIDADE}), data e
hora do fato, município, bairro, sexo, cor/raça, idade e escolaridade da
vítima, local do fato e instrumento utilizado. A versão utilizada neste
trabalho cobre registros até 30 de abril de 2026. Como o recurso é
atualizado periodicamente no portal, recomenda-se citar sempre a data de
acesso e a versão obtida, evitando divergências com revisões posteriores
do conjunto.

\subsection{Pré-processamento}
A base apresentou importantes desafios de qualidade, refletindo a realidade
de registros administrativos de segurança pública no Brasil. As principais
etapas de pré-processamento realizadas foram as seguintes.

\textbf{Tratamento de valores ausentes:} As variáveis \texttt{ESCOLARIDADE}
(74,15\% de valores faltantes) e \texttt{OCUPAÇÃO} (67,59\%) foram descartadas
por inviabilidade de imputação. A variável \texttt{COR/RAÇA} (5,12\% de
ausência) foi imputada pela moda --- a categoria ``Parda'' ---, enquanto
\texttt{IDADE} (2,11\%), \texttt{INSTRUMENTO} (1,87\%) e \texttt{LOCAL DO
FATO} (0,43\%) foram tratados por exclusão dos registros faltantes. Ressalte-se
que a imputação de cor/raça pela moda reforça a categoria majoritária e,
portanto, infla estatísticas descritivas baseadas nessa variável; tais
números devem ser lidos com essa ressalva (Seção~6).

\textbf{Engenharia de atributos temporais:} A partir da coluna de data, foram
criadas as variáveis: \texttt{ano}, \texttt{mes}, \texttt{dia\_semana} e
\texttt{periodo\_dia}, segmentado em quatro faixas: Madrugada (0h--5h),
Manhã (6h--11h), Tarde (12h--17h) e Noite (18h--23h).

\textbf{Discretização de idade:} A idade contínua foi transformada na variável
categórica \texttt{faixa\_etaria} com os grupos: 0--11, 12--17, 18--24,
25--29, 30--39, 40--59 e 60+.

\textbf{Agrupamento de variáveis categóricas:} O instrumento utilizado foi
consolidado em quatro categorias (\texttt{grupo\_instrumento}): Arma de fogo,
Arma branca, Espancamento e Outros. O local do fato foi consolidado em
\texttt{grupo\_local}: Espaço público, Entorno de casa, Ambiente interno e
Área externa/isolada.

\textbf{Normalização de nomes:} Nomes de municípios e bairros foram
normalizados (remoção de acentos, padronização de maiúsculas/minúsculas) para
viabilizar a junção com dados geográficos de centroide do IBGE.

Após o pré-processamento, a base resultante continha \textbf{20.369 registros}
(dos 21.217 originais) com todas as variáveis essenciais preenchidas,
utilizados de forma consistente em todas as análises descritivas e na
modelagem subsequentes.

\section{Estatística Descritiva e Análise Exploratória}
\subsection{Evolução Temporal}
A análise da evolução anual do CVLI em Alagoas revela uma tendência de queda
ao longo da série histórica, com pico em 2013 (2.181 registros) e mínimo em
2021 (1.028 registros, desconsiderando-se 2026, ainda parcial), conforme a
Figura~\ref{fig:serie_mensal}. Este
comportamento pode ser parcialmente associado a políticas públicas de
segurança implementadas no período, como o fortalecimento de programas
estaduais de combate à violência~\cite{ratton:2014}.

\begin{figure}[ht]
\centering
\includegraphics[width=.9\textwidth]{serie_mensal_al.png}
\caption{Evolução mensal do CVLI em Alagoas (2012--2026), evidenciando
  sazonalidade intra-anual e tendência de queda ao longo da série histórica.}
\label{fig:serie_mensal}
\end{figure}

\subsection{Padrões Temporais}
A distribuição por período do dia revela forte concentração noturna: o
período Noite (18h--23h) concentra 42,89\% das ocorrências, seguido por
Tarde (22,72\%), Madrugada (17,47\%) e Manhã (16,92\%). Combinando-se os
períodos Noite e Madrugada, \textbf{60,36\% dos CVLI ocorrem entre 18h e 5h}.
A análise por dia da semana indica que \textbf{37,88\% das ocorrências
concentram-se nos finais de semana} (sábado e domingo).

O mapa de calor da Figura~\ref{fig:heatmap} cruza o período do dia com o
dia da semana, revelando padrões claros e repetíveis: a concentração nas
noites de sexta-feira, sábado e domingo é consistente ao longo de toda a
série histórica, sugerindo uma sazonalidade semanal estrutural do fenômeno.

\begin{figure}[ht]
\centering
\includegraphics[width=.85\textwidth]{heatmap_dia_hora.png}
\caption{Mapa de calor: distribuição dos CVLI por dia da semana e período do
  dia (2012--2026). Cores mais intensas indicam maior concentração de
  ocorrências.}
\label{fig:heatmap}
\end{figure}

\subsection{Distribuição Geográfica}
A distribuição geográfica dos CVLI em Alagoas é marcadamente concentrada.
O município de Maceió, capital do estado, responde sozinho por
\textbf{33,33\% (6.789 registros)} de todas as ocorrências. As cinco cidades
com maior incidência — Maceió, Arapiraca (7,24\%), Rio Largo (4,15\%), União
dos Palmares (2,24\%) e Marechal Deodoro (2,16\%) — concentram \textbf{49,12\%}
dos CVLI do estado, ou seja, praticamente metade das ocorrências. A
Figura~\ref{fig:choropleth} apresenta a distribuição geográfica acumulada por
município por meio de um mapa de bolhas posicionado nos centroides
municipais (não há geometria de polígonos, de modo que não se trata de um
mapa coroplético propriamente dito).

\begin{figure}[ht]
\centering
\includegraphics[width=.6\textwidth]{choropleth_municipios.png}
\caption{Distribuição geográfica acumulada dos CVLI por município de Alagoas
  (2012--2026), representada por um mapa de bolhas posicionado nos centroides
  municipais. O tamanho e a cor dos marcadores indicam o volume de ocorrências,
  evidenciando a concentração em Maceió e em polos urbanos específicos do
  estado.}
\label{fig:choropleth}
\end{figure}

A persistência dessas manchas ao longo do tempo é ilustrada na
Figura~\ref{fig:janelas}, que compara a distribuição geográfica em quatro
janelas temporais agregadas: os mesmos polos urbanos concentram os casos de
forma recorrente em toda a série histórica.

\begin{figure}[ht]
\centering
\includegraphics[width=\textwidth]{evolucao_janelas.png}
\caption{Evolução espaço-temporal dos CVLI por janelas temporais em Alagoas
  (2012--2026). Cada painel representa um intervalo agregado de anos,
  permitindo comparar a persistência e a redistribuição das manchas criminais
  ao longo do tempo. O tamanho dos marcadores indica o volume de ocorrências
  por município em cada janela, evidenciando a concentração recorrente dos
  casos em polos urbanos específicos, especialmente na região de Maceió e em
  municípios de maior incidência no estado.}
\label{fig:janelas}
\end{figure}

\subsection{Perfil das Vítimas}
A análise do perfil demográfico das vítimas confirma padrões amplamente
documentados na literatura sobre violência letal no
Brasil~\cite{ipea:2021}: \textbf{93,44\% das vítimas são do sexo masculino};
\textbf{59,19\% têm entre 15 e 29 anos}; e \textbf{87,89\% são pardas ou
pretas} (na base limpa, em que a cor/raça ausente foi imputada como ``Parda'';
considerando apenas os registros com cor/raça informada, a proporção é
87,47\%). A interseção desses três fatores — homens jovens pardos ou pretos —
representa 50,59\% de todas as vítimas, evidenciando uma vulnerabilidade
estrutural acentuada nesse grupo.

Quanto ao instrumento utilizado, \textbf{79,07\% dos CVLI foram cometidos com
arma de fogo} (Projetil de Arma de Fogo -- PAF). Entre homens de 15 a 29
anos, essa proporção sobe para 85,27\%. Em contraste, arma branca ou
espancamento são relativamente mais frequentes entre vítimas do sexo feminino
(36,05\%) e em ambientes internos (27,24\%), sugerindo um perfil de violência
doméstica nesses casos.

Para verificar se os padrões descritos são estatisticamente significativos
e não artefatos de amostragem, aplicou-se o teste qui-quadrado de
independência. A associação entre o sexo da vítima e o instrumento utilizado
é altamente significativa ($\chi^2(3) = 551{,}2$; $p < 0{,}001$; V de
Cramér $= 0{,}165$): entre vítimas do sexo feminino, arma branca (25,2\%) e
espancamento (10,8\%) são proporcionalmente mais frequentes do que entre as
masculinas (13,0\% e 5,1\%), ao passo que a arma de fogo predomina entre
homens (80,6\% contra 56,9\%). Esse resultado sustenta estatisticamente a
hipótese de um perfil distinto de violência associado a vítimas femininas
(Seção~1). As concentrações temporais também rejeitam a hipótese de
independência --- período do dia $\times$ dia da semana ($\chi^2(18) =
350{,}1$; $p < 0{,}001$) e fim de semana $\times$ período do dia ($\chi^2(3)
= 107{,}0$; $p < 0{,}001$) --- embora com tamanho de efeito fraco ($V \approx
0{,}07$), indicando que o padrão é consistente, porém de magnitude modesta.

\section{Métodos Avaliados}
A análise foi conduzida em duas frentes complementares: (i)~análise
exploratória espaço-temporal para identificação de padrões e manchas
criminais; e (ii)~modelagem preditiva baseada em taxa de risco relativo por
perfil demográfico-contextual.

\subsection{Análise Espaço-Temporal}
A construção dos mapas de calor temporais utilizou tabelas dinâmicas
(\textit{pivot tables}) cruzando as variáveis \texttt{dia\_semana} e
\texttt{periodo\_dia}, agregando a contagem de ocorrências. Para a análise
geográfica, foram calculados centroides municipais a partir de dados do IBGE,
permitindo a construção de mapas de bolhas por centroide --- e não mapas
coropléticos, dada a ausência de geometria de polígonos --- e mapas de calor
interativos com a biblioteca \textit{Folium}~\cite{folium:2023}.

\subsection{Modelagem Preditiva de Risco}
A modelagem preditiva foi restrita ao município de Maceió por duas razões
objetivas: (i)~o município concentra 33,33\% de todos os registros de CVLI da
base, fornecendo volume suficiente para treinamento; e (ii)~dispõe de dados
populacionais por bairro no Censo IBGE 2022, o que permite calcular uma
\textbf{taxa de risco relativo por perfil demográfico}, mais justa do que a
simples contagem absoluta de ocorrências.

\textbf{Construção da matriz de cenários:} O modelo não foi treinado sobre
registros individuais, mas sobre uma matriz completa de cenários possíveis,
formada pela combinação de:

\begin{center}
$\text{bairro} \times \text{sexo} \times \text{faixa etária} \times
\text{dia da semana} \times \text{turno} \times \text{tipo de local}$
\end{center}

Para cada cenário foi contabilizado o número histórico de CVLI. Os cenários
sem ocorrência foram mantidos com valor zero, permitindo que o modelo
aprendesse também perfis de baixo risco.

\textbf{Normalização por população exposta:} Para evitar que bairros mais
populosos ou grupos majoritários dominassem a análise por volume bruto, foi
estimada a população de cada perfil a partir do Censo IBGE 2022:

\begin{equation}
\text{pop\_perfil} = \text{pop\_bairro}
  \times \text{prop\_sexo}
  \times \text{prop\_faixa\_etária}
\end{equation}

A variável-alvo do modelo foi a \textbf{taxa anualizada de CVLI por
100 mil habitantes}:

\begin{equation}
\text{taxa\_risco} = \frac{\text{crimes\_anuais}}{\text{pop\_perfil}}
  \times 100.000
\end{equation}

em que \texttt{crimes\_anuais} é a contagem histórica do cenário dividida
pelo número de anos da série. Na implementação atual utilizou-se o número de
anos-calendário distintos (15) como divisor; como o período efetivamente
coberto corresponde a aproximadamente 14,3 anos (jan/2012--abr/2026, com 2026
ainda parcial), esse divisor subestima levemente a taxa anualizada --- para
Maceió, por exemplo, a taxa real passa de 47,34 para cerca de 49,58/100 mil ---
e deve ser ajustado ao intervalo efetivo em versões futuras.

Essa taxa foi classificada em três níveis operacionais ---
\textit{baixo}, \textit{médio} e \textit{alto} --- com base em uma média de
referência recalibrada sobre os cenários com crimes reais, separada por sexo,
para evitar distorções causadas pelo grande número de combinações zeradas no
grid. Cabe registrar que essa média de referência é estimada sobre toda a
base, o que vincula a definição das classes ao conjunto completo (ver a
discussão sobre vazamento de alvo na Seção~6).

\textbf{Algoritmo e codificação:} O algoritmo escolhido foi o
\textbf{XGBoost Regressor com objetivo Poisson}~\cite{chen:2016}, adequado
para eventos raros baseados em contagens, pois penaliza corretamente a cauda
da distribuição sem inflar previsões em cenários de zero ocorrências. As
variáveis categóricas (bairro, sexo, faixa etária, dia da semana, turno e
local do fato) foram codificadas via \texttt{OneHotEncoder} antes do
treinamento, resultando em 6 atributos base expandidos por
\textit{one-hot encoding}.

\textbf{Explicabilidade com SHAP:} Para suporte à interpretação operacional,
o modelo incorpora valores SHAP (\textit{SHapley Additive
exPlanations})~\cite{lundberg:2017}, que quantificam a contribuição individual
de cada atributo na previsão de risco. Isso permite que gestores de segurança
compreendam, em linguagem natural, por que determinado perfil, bairro ou
horário recebeu alta classificação de risco.

É importante distinguir dois artefatos. O modelo \textbf{operacional},
efetivamente carregado pelo \textit{dashboard}, é o XGBoost Poisson Regressor
descrito acima, ajustado sobre todo o \textit{grid} de cenários sem conjunto de
teste apartado e sem validação cruzada. Separadamente, e apenas para fins de
comparação e linha de base, foi conduzido um \textbf{experimento de
classificação} sobre os registros individuais de Maceió (6.789 ocorrências),
em que cada ocorrência recebe um rótulo de risco (\textit{baixo}, \textit{médio}
ou \textit{alto}) derivado do percentil de frequência de seu contexto dentro do
grupo demográfico. Nesse experimento avaliaram-se \textbf{Regressão
Logística}~\cite{hastie:2009}, \textbf{Random Forest}~\cite{breiman:2001} e
\textbf{XGBoost} (classificador), com divisão treino/teste 80/20 estratificada
pela variável-alvo. Como o rótulo é construído a partir de contagens da base
inteira e essas mesmas contagens determinam as variáveis de entrada, esse
experimento sofre de \textbf{vazamento de alvo} --- suas métricas, reportadas na
Seção~6, são otimistas e não devem ser interpretadas como medida de
generalização (ver Seção~6).

\section{Métricas e Estratégia de Avaliação}
\subsection{Métricas}
A avaliação dos modelos utilizou as seguintes métricas: \textbf{Acurácia}
(proporção de classificações corretas no total); \textbf{F1-Score macro}
(média aritmética do F1-Score por classe, adequada para avaliar desempenho
equilibrado entre todas as classes); e \textbf{Precisão} e \textbf{Revocação}
por classe, para análise detalhada do comportamento do modelo em cada
categoria de risco. A Matriz de Confusão foi utilizada para inspeção dos
erros de classificação.

\subsection{Estratégia de Avaliação}
Além da divisão treino/teste (80/20), foi realizada \textbf{validação cruzada
estratificada com 5 \textit{folds}} para estimar a variância do desempenho e
garantir que os resultados não fossem sensíveis à partição específica dos
dados~\cite{hastie:2009}. Esta estratégia é recomendada para conjuntos de
dados de tamanho moderado, pois oferece estimativas mais robustas da
capacidade de generalização do modelo. A estratificação por classe garantiu
que a distribuição das categorias de risco fosse mantida em cada
\textit{fold}. Cabe a ressalva, no entanto, de que a validação cruzada acima
se aplica ao \textbf{experimento de classificação} e \textit{não} corrige o
vazamento de alvo: como o rótulo de risco é definido a partir de agregados de
toda a base, a separação em \textit{folds} não isola informação entre treino e
teste, de modo que as estimativas permanecem otimistas (Seção~6). O modelo
Poisson operacional, por sua vez, não foi submetido a partição de teste nem a
validação cruzada no \textit{pipeline} de treino.

\section{Resultados e Discussão}
\subsection{Desempenho dos Modelos (Experimento de Classificação)}
A Tabela~\ref{tab:modelos} apresenta os resultados do experimento de
classificação descrito na Seção~4.2 --- treinado sobre os registros
individuais de Maceió, distinto do modelo Poisson operacional --- no conjunto
de teste (20\% dos dados). \textbf{Estes resultados devem ser lidos com forte
ressalva}, pois o rótulo de risco sofre vazamento de alvo, conforme detalhado
adiante.

\begin{table}[ht]
\centering
\caption{Comparação de desempenho no experimento de classificação sobre
  registros individuais de Maceió (conjunto de teste). \textbf{Atenção:}
  métricas otimistas devido a vazamento de alvo (ver Seção~6); não
  correspondem ao modelo Poisson operacional.}
\label{tab:modelos}
\begin{tabular}{|l|c|c|}
\hline
\textbf{Modelo} & \textbf{Acurácia} & \textbf{F1-Score Macro} \\
\hline
Regressão Logística  & 84,02\% & 83,98\% \\
Random Forest        & 89,40\% & 89,19\% \\
\textbf{XGBoost}     & \textbf{92,42\%} & \textbf{92,35\%} \\
\hline
\end{tabular}
\end{table}

Nesse experimento, o XGBoost obteve os maiores valores em ambas as métricas,
superando a Regressão Logística em 8,4 pontos percentuais e o Random Forest em
3,0 pontos percentuais de acurácia, com 94,70\% ± 0,71\% na validação cruzada
de 5 \textit{folds}.

\textbf{Ressalva metodológica (vazamento de alvo).} As métricas acima
\emph{superestimam} a capacidade real de generalização e não devem ser
interpretadas como evidência de um preditor de alta acurácia. O rótulo
\texttt{classe\_risco} é construído a partir de contagens de frequência
calculadas sobre \emph{toda} a base, antes da divisão treino/teste, e essas
mesmas contagens são determinadas pelas variáveis usadas como entrada do
modelo. Em consequência, o alvo é praticamente uma função determinística das
\textit{features}: ao reproduzir a partição, verificou-se que \textbf{80,78\%
das linhas de teste compartilham com o treino a mesma combinação que define o
rótulo}, de modo que o modelo essencialmente recupera um valor já presente nos
dados de treino. A validação cruzada \emph{não} mitiga esse problema, pois o
vazamento ocorre na própria definição do alvo. Em uma verificação fora da
amostra do modelo Poisson operacional (divisão 80/20 sobre combinações de
cenário não vistas), o poder preditivo é fraco ($R^2 \approx 0{,}10$), o que
reforça que esse componente deve ser usado como ferramenta de
pontuação/triagem de risco relativo, e não como classificador de alta
acurácia.

A Figura~\ref{fig:comparacao} apresenta a comparação visual entre os modelos,
e a Tabela~\ref{tab:report} detalha as métricas por classe para o modelo
XGBoost.

\begin{figure}[H]
\centering
\includegraphics[width=.8\textwidth]{comparacao_modelos.png}
\caption{Comparação de desempenho entre os modelos avaliados: Regressão
  Logística, Random Forest e XGBoost.}
\label{fig:comparacao}
\end{figure}

\begin{table}[ht]
\centering
\caption{Métricas por classe --- XGBoost no experimento de classificação.
  Valores afetados pelo vazamento de alvo discutido na Seção~6.}
\label{tab:report}
\begin{tabular}{|l|c|c|c|}
\hline
\textbf{Classe} & \textbf{Precisão} & \textbf{Revocação} & \textbf{F1-Score} \\
\hline
Alta   & 94\% & 96\% & 95\% \\
Baixa  & 93\% & 94\% & 94\% \\
Média  & 90\% & 87\% & 88\% \\
\hline
\end{tabular}
\end{table}

A maior taxa de confusão ocorreu entre as classes \textit{média} e as demais,
o que é esperado dado o caráter limítrofe dos casos de risco intermediário.
A classe \textit{alta} obteve o melhor desempenho (F1 = 95\%), o que é
particularmente relevante do ponto de vista operacional, pois erros de
classificação na direção de subestimar o risco alto têm maior custo social.

\subsection{Insight Acionável}
Os resultados da análise espaço-temporal convergem para uma recomendação
central: \textit{60,36\% dos CVLI em Alagoas ocorrem no período noturno e de
madrugada (18h--5h), concentrados em manchas criminais identificáveis em
Maceió e em outras quatro cidades que, juntas, respondem por praticamente
metade (49,12\%) de todos os CVLI do estado. Esse padrão sugere que o
patrulhamento uniforme aloca parte expressiva dos recursos em zonas e períodos
de baixo risco --- a magnitude exata desse desperdício, porém, não foi
quantificada neste estudo e permanece como hipótese a ser testada. Diante
disso, recomendamos avaliar um protocolo de patrulhamento preditivo dinâmico,
alocando viaturas por faixa de horário e microrregião conforme as
concentrações identificadas. A recomendação apoia-se sobretudo na
concentração espaço-temporal evidenciada pela análise descritiva e pelo
mapa de calor (60,36\% dos CVLI no período noturno e nas cinco principais
microrregiões); ela \emph{não} se apoia nas métricas do experimento de
classificação, que são otimistas por vazamento de alvo (Seção~6). Após seis
meses, o efeito da ação poderia ser avaliado pela redução do CVLI nas áreas de
intervenção, acompanhando-se a taxa de incidentes nas zonas adjacentes como
controle do efeito de deslocamento. A título de meta operacional --- e não de
previsão derivada do modelo ---, a literatura de \textit{hot spots policing}
sugere quedas da ordem de 10\% a 20\% em áreas de intervenção concentrada,
parâmetro que poderia balizar a avaliação.}

Esta recomendação é consistente com a literatura de \textit{hot spots
policing}, que demonstra reduções médias de 10--20\% na criminalidade em
áreas de intervenção concentrada~\cite{sherman:1989, weisburd:2016}.

\subsection{Limitações}
Algumas limitações devem ser consideradas na interpretação dos resultados.
Primeiro, e de forma mais crítica, o \textbf{experimento de classificação}
(Tabela~\ref{tab:modelos}) sofre de vazamento de alvo: o rótulo de risco é
derivado de contagens calculadas sobre toda a base e essas mesmas contagens
determinam as \textit{features}, de modo que 80,78\% das linhas de teste
reaparecem, na prática, no treino. Suas métricas são, portanto, otimistas e não
medem generalização. O modelo Poisson operacional, por sua vez, foi ajustado
sobre todo o \textit{grid} sem conjunto de teste apartado; uma verificação
posterior fora da amostra indicou poder preditivo fraco ($R^2 \approx 0{,}10$),
o que o caracteriza como ferramenta de triagem de risco relativo, não como
preditor de alta acurácia.

Segundo, a base apresenta altas taxas de valores ausentes nas variáveis de
escolaridade e ocupação, limitando análises socioeconômicas mais aprofundadas;
além disso, a cor/raça foi imputada pela moda (``Parda''), o que infla
estatísticas descritivas baseadas nessa variável e deve ser considerado ao
ler a proporção de 87,89\% de vítimas pardas ou pretas. Terceiro, o modelo
preditivo foi treinado exclusivamente com dados de Maceió, podendo não
generalizar para municípios menores com dinâmicas distintas. Quarto, a ausência
de coordenadas geográficas precisas para cada ocorrência limitou a análise ao
nível de município e bairro; as visualizações ``geográficas'' são mapas de
bolhas posicionados em centroides (no caso dos bairros, coordenadas
aproximadas geradas por deslocamento determinístico em torno do centro da
cidade), e não mapas coropléticos nem produto de clusterização espacial.
Quinto, a análise descritiva foi complementada por testes qui-quadrado de
independência, que confirmam a significância estatística das associações
identificadas (Seção~3.4); contudo, não foi conduzida análise causal, de modo
que afirmações sobre o efeito de intervenções permanecem hipóteses. Sexto, a
taxa anualizada utiliza como divisor o número de anos-calendário (15) em vez
do intervalo efetivo (cerca de 14,3 anos), subestimando levemente as taxas.
Por fim, o modelo não incorpora \texttt{COR/RAÇA} como atributo de
normalização da taxa de risco, pois dados populacionais por raça/cor na
granularidade de bairro não estão disponíveis no Censo IBGE 2022.

\section{Conclusão}
Este trabalho mostrou que os Crimes Violentos Letais Intencionais em
Alagoas apresentam padrões espaço-temporais claros e consistentes. A
concentração de 60,36\% das ocorrências no período noturno e de madrugada,
de 37,88\% nos finais de semana e de 33,33\% no município de Maceió evidencia
que o fenômeno da violência letal não é aleatório — é estruturado e recorrente
ao longo da série histórica.

A partir desses padrões, construiu-se um modelo XGBoost Poisson, treinado sobre
uma matriz de cenários normalizada por estimativas populacionais do Censo IBGE
2022, como ferramenta operacional de pontuação de risco relativo por perfil.
Cabe registrar que a verificação fora da amostra indica poder preditivo ainda
limitado ($R^2 \approx 0{,}10$) e que as métricas elevadas do experimento de
classificação (92,42\% de acurácia; 94,70\% em validação cruzada) refletem
vazamento de alvo, não devendo ser tomadas como medida de acurácia do
instrumento. Ainda assim, o trabalho demonstra ser viável organizar dados
públicos já disponíveis em um instrumento de apoio à decisão, sem custo
adicional de coleta, desde que aprimorados o desenho do alvo e a validação do
modelo. A título de referência, a literatura de \textit{hot spots policing}
associa intervenções concentradas a reduções da ordem de 10\% a 20\% nas
áreas-alvo, parâmetro que pode orientar a avaliação de um eventual protocolo de
patrulhamento preditivo.

Como trabalhos futuros, destacam-se: (i)~expansão do modelo para os 102
municípios de Alagoas; (ii)~integração com dados socioeconômicos do IBGE e
com variáveis de raça/cor para análise de fatores estruturais; (iii)~implantação
de um \textit{dashboard} interativo para uso operacional pela SSP-AL com
explicações em linguagem natural via SHAP; e (iv)~aplicação de modelos de
séries temporais para previsão de tendências de médio prazo.

A base de dados utilizada é pública e está disponível em:
\url{dados.al.gov.br/catalogo/dataset/cvli-2012-a-2023-base-microdados}.

\bibliographystyle{sbc}
\bibliography{sbc-template}
\end{document}